\section{BEVEZETÉS}

Az autóutak burkolatának megléte és minősége rendkívül fontos a mai
világban, ahol nap mint nap rengeteg jármű használja ezeket az utakat
közlekedésre. A használatból és környezeti tényezőktől (esővíz
beszivárgás, fagyás, olvadás) a megépítés és vagy karbantartás nem
megfelelő minőségéből fakadóan úthibák képződhetnek és az idő múlásával
ezek egyre súlyosodhatnak, ha nincsnek időben felmérve és megfelelően
javítva \cite{abed_analysis_2023} .Ezt a jelenséget nevezzük kátyúnak, ami az útburkolat
felszínen található durva, egyenletlen tál alakú mélyedéshez
hasonlítható, amelynek szélessége és mélysége is széles skálán mozog
\cite{miller_distress_2014}, \cite{safyari_review_2024}.

A kátyúk a közúti infrastruktúra olyan hibái, amelyek világszerte
jelentős, problémákat okoznak. Talán a leginkább szembetűnő probléma az
utazási menetkomfort romlása amikor a jármű belehajt egy kátyúba.
Továbbá közlekedésbiztonsági veszélyt hordoznak magukban akár a jármű
feletti tejes uralom elvesztéséhez vezethetnek, különösen nagy
sebességnél, főként a kiszolgáltatott járművek (például kerékpárosok és
motorosok) körében ezzel növelve a balesetek bekövetkezési kockázatát
\cite{yalew_impact_2023}. Gazdasági szempontból a kátyúk jelentős pénzügyi terhet
jelentenek az úthálózat fenntartó szervezeteknek és az utat használó
közlekedők számára egyaránt. Míg a közlekedő járművek esetén jelentősen
befolyásolják és növelhetik a járművek karbantartási költségeit azzal,
hogy csökkentik az abroncsok, futómű alkatrészek és egyéb mechanikus
alkatrész élettartamát. Addig a kátyúk detektálása és javítása hatalmas
összegeket emészthet fel kormányzati oldalról nézve. Az Egyesült
Királyság kormánya olyan közleményt adott ki, amelyben azt írták, hogy
2020 és 2025 között időszakban több mint 5 milliárd fontot terveznek
kátyú detektálásra és javításra fordítani \cite{safyari_review_2024}. „Ezenkívül a Világbank
szerint az optimális állapotban lévő utak mennyisége egy ország
gazdasági erejének és versenyképességének mutatójaként használható.''
\cite{martinez-rios_review_2022}

Az előbb felsoroltak miatt a kátyúkkal kapcsolatos kutatások a mai napig
aktívak, és engem is arra sarkallott, hogy érdemes ezzel a témával
foglalkozni. Illetve van még egy tényező, amikor autót vezetek gyakran
találom magam szembe kátyúkkal az utakon ezeket a lehetőségekhez mérten
rendre próbálok kikerülni úgy, hogy ezzel másokat ne veszélyeztessek az
úton, így próbálom óvni az autó szerkezeti állapotát és a saját magam
vagy adott esetben az utasaim menetkomfortját. Bármennyire is igyekszem
figyelni vezetés közben az utat nyilvánvalóan előfordul, hogy későn,
vagy egyáltalán észre se veszem a kátyút az úton és akarva akaratlanául
áthajtok rajta.

Meglátásom szerint kátyúk mindig is lesznek amig javításuk ennyire drága
és körülményes. Ebbe beletörődve és autóval az úton töltött
tapasztalataimat figyelembe véve jutottam arra, hogy valós időben
szeretném a kátyúkat megbízhatóan felismerni, majd azok minél pontosabb
elhelyezkedését megállapítani az úton belül.